\documentclass[11pt]{article}

\usepackage[T1]{fontenc}
\usepackage[utf8]{inputenc}
\usepackage{lmodern}
\usepackage[margin=1in]{geometry}
\usepackage{microtype}
\usepackage{booktabs}
\usepackage{array}
\usepackage{graphicx}
\usepackage[font=small,skip=6pt]{caption}
\usepackage{longtable}
\usepackage{xcolor}
\usepackage[numbers,sort&compress]{natbib}
\usepackage{url}
\usepackage[hidelinks]{hyperref}
\usepackage{enumitem}
\setlist{nosep,leftmargin=*}

\newcommand{\ceoldef}[2]{\expandafter\def\csname ceol:#1\endcsname{#2}}
\newcommand{\cv}[1]{\ifcsname ceol:#1\endcsname\csname ceol:#1\endcsname\else
  \textbf{\textcolor{red}{??\detokenize{#1}??}}\PackageWarning{ceol}{Undefined value \detokenize{#1}}\fi}
\ceoldef{acc_mean_1}{278,722}
\ceoldef{acc_mean_8}{1,089,968}
\ceoldef{alone_enable}{78}
\ceoldef{alone_wiktionary_pct}{43.7\%}
\ceoldef{apos_both_flags}{41}
\ceoldef{apos_neither}{2,033}
\ceoldef{apos_rows}{140,851}
\ceoldef{arch_eol_universe}{9,653,999}
\ceoldef{arch_merge_total}{9,867,866}
\ceoldef{archival_eol_revision}{c1139698}
\ceoldef{audit_beyond_frame_added}{33}
\ceoldef{audit_frame_size}{9,787,841}
\ceoldef{audit_in_frame_added}{133,842}
\ceoldef{audit_misses}{140,974}
\ceoldef{audit_not_shipped}{5,400}
\ceoldef{audit_wordfreq_left_out}{7,099}
\ceoldef{c4_additions}{133,875}
\ceoldef{census_candidates}{213,867}
\ceoldef{census_not_shipped}{26,601}
\ceoldef{channels_shown}{12}
\ceoldef{char_rows_apostrophe}{140,851}
\ceoldef{char_rows_hyphen}{63,421}
\ceoldef{char_rows_nonbase_max}{2,794}
\ceoldef{char_rows_nonbase_min}{39}
\ceoldef{char_rows_space}{294,865}
\ceoldef{chars_total}{50}
\ceoldef{class_combos_nonplain}{47}
\ceoldef{co_both}{127,590}
\ceoldef{co_gb_only}{7,900,078}
\ceoldef{co_wf_only}{5,313}
\ceoldef{code_commit}{b03add02b151}
\ceoldef{combos_shown}{12}
\ceoldef{combos_total}{32}
\ceoldef{config_count}{7}
\ceoldef{cont_hunspell_us_in_hunspell_ca}{98.9\%}
\ceoldef{cont_scowl_in_enable}{39.2\%}
\ceoldef{cont_wiktionary_in_scowl}{43.1\%}
\ceoldef{cont_wiktionary_in_wordnet}{8.1\%}
\ceoldef{cont_wordnet_in_scowl}{94.5\%}
\ceoldef{cont_wordnet_in_wiktionary}{83.8\%}
\ceoldef{controls_checked}{20}
\ceoldef{core_layer_eol_invalid}{95,245}
\ceoldef{core_layer_eol_valid}{202,423}
\ceoldef{core_pct}{4.2\%}
\ceoldef{core_rows}{417,864}
\ceoldef{cov_hunspell_us_b5plus}{99.9\%}
\ceoldef{cov_wordnet_b5plus}{82.7\%}
\ceoldef{covw_wordnet_b5plus}{57.4\%}
\ceoldef{enable_in_scowl}{99.7\%}
\ceoldef{eol_revision}{c550751e}
\ceoldef{eol_universe_rows}{9,653,667}
\ceoldef{family_count}{5}
\ceoldef{filter_count}{5}
\ceoldef{flag_any}{482,131}
\ceoldef{flag_contraction}{1,841}
\ceoldef{flag_plain}{9,782,142}
\ceoldef{flag_plain_pct}{95.3\%}
\ceoldef{flag_s_final}{137,018}
\ceoldef{freq3_unlisted}{378}
\ceoldef{gb_listed_us_blind}{9,147}
\ceoldef{gb_not_us}{146,322}
\ceoldef{gb_p50_corpus}{2.51}
\ceoldef{gb_p50_one}{2.96}
\ceoldef{gb_p50_two_plus}{3.95}
\ceoldef{gb_sole_pct_all}{77.0\%}
\ceoldef{gb_with_corpus_pct}{99.9\%}
\ceoldef{gb_with_one_pct}{56.9\%}
\ceoldef{gb_with_two_plus_pct}{84.1\%}
\ceoldef{generated_pct}{3.22\%}
\ceoldef{generated_rows}{330,363}
\ceoldef{grammar_accented}{19}
\ceoldef{grammar_base}{26}
\ceoldef{grammar_ligatures}{2}
\ceoldef{hb_tier2}{27,458}
\ceoldef{hb_ws}{2,558}
\ceoldef{hb_ws_gb}{1,978}
\ceoldef{hb_ws_wordnet}{588}
\ceoldef{hb_ws_wordnet_gb}{8}
\ceoldef{held_back_rows}{30,016}
\ceoldef{hun_all4_hold_listed}{268,042}
\ceoldef{hun_all4_hold_some_pct}{60.3\%}
\ceoldef{hun_any_some_pct}{61.6\%}
\ceoldef{hunspell_dictionaries}{4}
\ceoldef{hyph_wf_false}{53,840}
\ceoldef{inh_unlisted_gb}{8,024,365}
\ceoldef{inh_unlisted_gb_only}{7,900,077}
\ceoldef{inh_unlisted_gb_pct_all}{78.18\%}
\ceoldef{inh_unlisted_gb_pct_single}{80.49\%}
\ceoldef{leak_applicable}{12}
\ceoldef{leak_auc_ho_max}{0.7752}
\ceoldef{leak_auc_ho_min}{0.6170}
\ceoldef{leak_auc_max}{0.7823}
\ceoldef{leak_auc_min}{0.6044}
\ceoldef{leak_strata_max}{7}
\ceoldef{leak_strata_min}{4}
\ceoldef{len_corpus_p50}{8}
\ceoldef{len_multiword_p50}{14}
\ceoldef{len_none_p50}{10}
\ceoldef{len_one_p50}{10}
\ceoldef{len_two_plus_p50}{9}
\ceoldef{limit_gb_all}{8,761,597}
\ceoldef{limit_no_panel_no_corpus}{846,459}
\ceoldef{list_enable_blind}{187,266}
\ceoldef{list_hunspell_us_blind}{61,747}
\ceoldef{list_wiktionary_listed}{831,973}
\ceoldef{list_wordnet_listed}{80,441}
\ceoldef{listed_no_corpus}{347,609}
\ceoldef{listed_no_corpus_pct}{31.9\%}
\ceoldef{listed_single}{1,089,968}
\ceoldef{m5_k0}{8,872,182}
\ceoldef{m5_k0_pct}{90.7\%}
\ceoldef{m5_k1_pct_of_listed}{59.4\%}
\ceoldef{m5_k5_pct_of_listed}{5.3\%}
\ceoldef{m5_listed}{909,960}
\ceoldef{model_basis_entries}{344,830}
\ceoldef{mw_both}{19,511}
\ceoldef{mw_wiktionary_only}{231,382}
\ceoldef{mw_wordnet}{63,483}
\ceoldef{mw_wordnet_only}{43,972}
\ceoldef{one_family_pct_listed}{61.7\%}
\ceoldef{p1_c1_exceptions}{187,266}
\ceoldef{p1_c4_additions}{128,475}
\ceoldef{p1_eol_invalid}{9,308,837}
\ceoldef{p1_eol_valid}{344,830}
\ceoldef{p1_spaced_c015}{294,865}
\ceoldef{p2_listed}{3,953}
\ceoldef{p2_listed_pct}{1.2\%}
\ceoldef{p2_none}{219}
\ceoldef{p4_channels}{26}
\ceoldef{p4_top_listed_pct}{5.7\%}
\ceoldef{p4_top_rows}{9,009,717}
\ceoldef{pairs_both_no_common}{5,660}
\ceoldef{pairs_none}{4,291}
\ceoldef{pairs_one}{72,135}
\ceoldef{pairs_share}{93,019}
\ceoldef{pairs_total}{175,105}
\ceoldef{panel_set_count}{13}
\ceoldef{pin_spaced}{294,865}
\ceoldef{prov_channel_and_containment}{9,417,814}
\ceoldef{prov_channel_only}{846,459}
\ceoldef{prov_containment_only}{0}
\ceoldef{prov_neither}{0}
\ceoldef{published_set_count}{8}
\ceoldef{reach_threshold}{99\%}
\ceoldef{repin_left}{332}
\ceoldef{resources_count}{26}
\ceoldef{restricted_words_checked}{1,985}
\ceoldef{rows_all}{10,264,273}
\ceoldef{rows_eol_valid}{344,830}
\ceoldef{rows_exceptions}{209,883}
\ceoldef{rows_features}{9,969,408}
\ceoldef{rows_listed}{1,089,968}
\ceoldef{rows_listed_pct}{10.9\%}
\ceoldef{rows_panel_listed}{1,384,833}
\ceoldef{rows_pin}{10,296,274}
\ceoldef{rows_restricted}{1,985}
\ceoldef{rows_sample}{39,938}
\ceoldef{rows_shippable}{10,294,289}
\ceoldef{rows_single}{9,969,408}
\ceoldef{rows_spaced}{294,865}
\ceoldef{rows_universe}{10,264,273}
\ceoldef{s3_rows}{9,782,142}
\ceoldef{s_final_token_only}{11,641}
\ceoldef{sample_prefix}{00}
\ceoldef{span_p50_corpus}{181}
\ceoldef{span_p50_one}{160}
\ceoldef{span_p50_two_plus}{252}
\ceoldef{strict2_core}{366,934}
\ceoldef{strict_core}{372,809}
\ceoldef{strict_core_less}{45,055}
\ceoldef{strict_core_pct}{3.7\%}
\ceoldef{tier_corpus_pct}{80.6\%}
\ceoldef{tier_corpus_rows}{8,032,981}
\ceoldef{tier_count}{4}
\ceoldef{u5_gb_p50_k0}{2.51}
\ceoldef{u5_gb_p50_k5}{5.89}
\ceoldef{u5_rows_k5}{48,408}
\ceoldef{us_and_gb}{106,944}
\ceoldef{us_not_gb}{14,308}
\ceoldef{v1_rows}{10,001,409}
\ceoldef{var_groups}{191,985}
\ceoldef{var_rows}{402,875}
\ceoldef{variant_largest}{7}
\ceoldef{variant_patterns}{16}
\ceoldef{variant_patterns_shown}{10}
\ceoldef{version}{v1.0}
\ceoldef{wf_audit_rows}{10,191}
\ceoldef{wf_census_rows}{21,587}
\ceoldef{wf_sole_rows}{7,389}
\ceoldef{wikt_alone_m5}{426,798}
\ceoldef{wikt_alone_m5_pct}{46.9\%}
\ceoldef{withheld_feature_count}{12}
\ceoldef{withheld_list_count}{5}
\ceoldef{withheld_model_feature_count}{2}
\ceoldef{withheld_reference_feature_count}{10}
\ceoldef{zipf_p50_corpus}{1.35}
\ceoldef{zipf_p50_one}{1.60}
\ceoldef{zipf_p50_two_plus}{2.05}
\ceoldef{zipf_with_corpus_pct}{1.7\%}
\ceoldef{zipf_with_one_pct}{6.7\%}
\ceoldef{zipf_with_two_plus_pct}{30.1\%}


\DeclareUrlCommand\col{\urlstyle{tt}}
\renewcommand{\UrlBreaks}{\do\_\do\.\do\-\do\/\do\:\do\@}
\newcommand{\resource}[1]{\par\medskip\noindent\textbf{#1}\par}
\newcommand{\resfield}[2]{\noindent\hangindent=1.5em\hangafter=1\emph{#1:} #2\par}

\title{The CEOL Universe: an inventory of candidate English strings and the evidence for each}
\author{Ryan Kamp}
\date{November 26, 2026}

\begin{document}
\maketitle

\begin{abstract}
Studies that need to know whether a string already exists have screened it against
references of their own choosing and used the result to exclude items. The CEOL Universe is an
inventory of \cv{rows_all} candidate English strings that records the answer of each source instead
of a verdict. For each single-token row, it records which of \cv{published_set_count} published word
lists list the string, which of them cannot hold strings of its spelling class, and whether two
corpora contain it; for every row, it records how the string entered the pool. No row says that a
string is a word. Of the \cv{rows_single} single-token rows, \cv{rows_listed} (\cv{rows_listed_pct})
are listed by at least one list, and \cv{tier_corpus_rows} (\cv{tier_corpus_pct}) are listed by
none and held by a corpus. For \cv{gb_sole_pct_all} of all rows, the Google Books 1-grams are the only
published containment. Of the listed single-token rows, \cv{one_family_pct_listed} are listed by
exactly one panel family, and \cv{alone_wiktionary_pct} by Wiktionary alone. A reader with the data
derives a word list from the inventory under a definition they state, with code that reproduces it.
\end{abstract}

\section{Introduction}

Word lists disagree. A spell-checker's dictionary, a word-game list, a lexical database, and a
corpus's vocabulary each draw the line around English in a different place, and none is wrong by
its own rule. Studies that need to know whether a string ``exists'' have built their own checks:
the three benchmarks of invented terms in Section~\ref{sec:related} each used a different screening
procedure and used the result to exclude items.

The CEOL Universe takes the other approach. It keeps the candidate strings its sources gave that its
admission rules accept, under one executable spelling grammar, and it records each source's answer
beside the string, dated by the source's pinned version. A reader who wants a word list states a
definition (for example, ``listed by two or more panel families, single tokens only, and no string
the English OpenList's synthetic engine made'') and, with the data, obtains a list, a count, and the
code that reproduces both.

This paper describes version~\cv{version}. Its contributions are:
\begin{itemize}
\item \textbf{The inventory:} \cv{rows_all} rows, each a candidate string with an identifier that
  depends on the string alone, so the same string has the same identifier in every version that
  holds it.
\item \textbf{The evidence semantics:} every single-token row has a status in every published list:
  1 (listed; for SCOWL, a positive size tier instead), 0 (not listed), or $-1$ (the list cannot hold
  strings of this class, so it was not looked up). On a multiword row, the evidence columns are null (not
  measured). Directed containment and coverage by frequency are computed over the rows a list can
  hold.
\item \textbf{A comparison of the published lists over the same strings} (Section~\ref{sec:compared}).
  From those statuses, the paper states how far the lists contain one another, what each family
  holds alone, and how listing depends on frequency.
\item \textbf{A record of how each row came in} (Sections~\ref{sec:construction}
  and~\ref{sec:contents}), kept apart from the evidence for it.
\end{itemize}

\paragraph{What the universe is not.} It is not a dictionary: it has no definitions, senses, or
pronunciations. It does not decide whether a string is a word, and it never calls a candidate
valid. It is not a sample of English text, and its counts are counts of strings, not of uses.

\section{Related work}
\label{sec:related}

Appendix~\ref{app:resources} lists, for \cv{resources_count} English word lists, lexical databases, and
frequency resources, what each is, its stated size and license, and what it records per entry, as
read on October 26, 2026 unless an entry or its reference gives another date; some have changed since.

\paragraph{Records across resources.} Most word lists record only their own membership. A related
multi-resource record is UBY~\cite{gurevych2012uby}, which integrates nine lexical-semantic
resources, English WordNet and English Wiktionary among them, in ten lexicons, and counts lemmas,
within a language and selected parts of speech, by how many lexicons hold them. The cited paper describes no pinned file versions, no value for a
resource that cannot hold a string or did not measure it, no word lists, and no spelling grammar. Among the word lists in
Appendix~\ref{app:resources}, the nearest to a per-source record are ESDB, whose tags sometimes name
the source list a word came from~\cite{atkinson2026a}, and Wordnik's paid game dataset, which labels
words by source~\cite{wordnik2026}. The 12dicts lists are an earlier agreement design: several of its
classic lists keep the words that a stated number of its 12 source dictionaries hold, with
list-specific rules and exceptions~\cite{beale2026}.
Michel and colleagues measured dictionary coverage by frequency: sampling bands of the Google Books
lexicon, they found that the Oxford English Dictionary and Merriam-Webster's Unabridged covered
high-frequency words well and low-frequency words much less~\cite{michel2011}.

\paragraph{Documentation standards.} Datasheets for Datasets~\cite{gebru2021}, Data
Statements~\cite{bender2018,bender2021,mcmillanmajor2024}, Data Cards~\cite{pushkarna2022}, the
Hugging Face dataset card template~\cite{huggingface2026}, Croissant and its Responsible AI
extension~\cite{akhtar2024,mlcommons2026,mlcommons2024,akhtar2024a}, and the FAIR
principles~\cite{wilkinson2016} each set out what a dataset's documentation should cover, and
they differ in what they ask. Lexicographic resources have their own standards: TEI
Lex-0~\cite{teilex2026}, the ELEXIS Seal of Compliance~\cite{tasovac2021}, the European Language
Grid's metadata~\cite{labropoulou2020}, and the LRE Map~\cite{calzolari2012}. The Research Data
Alliance's recommendations for evolving data ask for versioned, timestamped data and a checksum per
query result~\cite{rauber2015}. In the versions cited, none of these dataset standards has an
explicit provision for a canary or for stating a dataset's exposure to contamination. BetterBench, a framework for assessing benchmarks, does: it
asks benchmark developers to include a globally unique identifier or canary string in their public
evaluation code and data files, or to encrypt the data and publish the key, and it notes that developers cannot enforce a
canary's use~\cite{reuel2024betterbench}.

\paragraph{Contamination and canaries.} Sainz and colleagues separate guideline, raw-text, and
annotation contamination~\cite{sainz2023}, and surveys organize the detection methods by
technique~\cite{deng2024,xu2024}. BIG-bench puts a GUID canary in every task definition file so that its
data can be filtered out of web-scraped corpora and a model can be tested afterward for training on
it~\cite{srivastava2023}; GPQA and BrowseComp embed canaries that contain the BIG-bench
string~\cite{rein2024,wei2025}. As a filter, a canary helps only trainers who know of it and filter
on it, and it does nothing against a developer who means to use the data~\cite{chen2025}; third-party copies of
benchmarks have been found with the canary removed~\cite{han2025}. A memorization canary, as in the
Secret Sharer study, is a different object: random sequences inserted into training data to measure
their exposure~\cite{carlini2019}. Several benchmarks make contamination questions answerable by
date. LiveCodeBench tags each problem with its release date, so a model can be scored on problems
released after its training cutoff, and uses the model's release date where the cutoff is
unknown~\cite{jain2025}. Deng and colleagues caution that a time cutoff works only if data before and
after it differ~\cite{deng2024}.

\paragraph{Existence checks.} Benchmarks of invented terms check existence while they are built or
analyzed, and they drop what the check finds. PhantomBench keeps a concept only if it has zero exact
matches in the Dolma v1.7 corpus~\cite{jung2026}. HypoTermQA removed, at its search stage, every
generated term for which a quoted web search returned a nonzero result count~\cite{uluoglakci2024}. A vocabulary-test benchmark
excluded two Spanish nonwords from its analysis after finding them in dictionaries~\cite{martinez2024}.
KoAGE builds a dated, layered exclusion lexicon from a dictionary, Korean Wiktionary, and dated corpus
and neologism layers, and uses it as a filter for collecting Korean neologisms~\cite{kim2026koage}.
The universe records, per string, which published lists hold it and which were not measured, as a
stored record rather than a filter.

\section{Design}
\label{sec:design}

\subsection{The candidate and the grammar}

A row of the universe is a \emph{candidate}: a string that passes the spelling grammar, with the
record of where it came from. The grammar (\col{pipeline/ceol/grammar.py}) is executable. It puts a
source entry in Unicode normal form C, replaces the typographic apostrophe with the ASCII apostrophe
and the no-break space with a space, lowers its case, and puts it in normal form C again; it then
rejects the entry if any character outside its charset remains.
It admits the result when every character is in a closed charset of \cv{chars_total}: the
\cv{grammar_base} letters a to z, \cv{grammar_accented} accented letters, \cv{grammar_ligatures}
ligatures, the hyphen, the apostrophe, and the space. It also sets rules of form: spaces are single
and interior; hyphens are interior, never doubled, and every hyphen-separated segment carries a
letter; an apostrophe is interior or leading, never final or doubled; every token carries a letter;
and every space-separated token meets the single-token rules. A space separates tokens, so a multiword term is one candidate.

The grammar marks each candidate with its spelling classes: \col{accented} (an accent other than
the diaeresis), \col{apostrophe-form}, \col{contraction} (an apostrophe form the grammar's rule takes
for a contraction), \col{diaeresis}, \col{hyphenated}, \col{ligature}, \col{s-final} (a token ending in
an apostrophe and s), and \col{spaced}. A candidate can
carry several; one with none is \emph{plain}.

\subsection{The evidence semantics}

For each published list, a single-token row holds 1 where the list lists the string, 0 where it does
not, and $-1$ where the list is structurally blind to the row's class: a list built without hyphens
cannot hold a hyphenated string. For SCOWL, the value is the size of the smallest SCOWL list that
holds the string, 0 where none does, and $-1$ where the class does not apply.

The evidence columns are null on every multiword row. Two published lists hold multiword terms,
WordNet and Wiktionary, and a multiword row's listing by them is recorded in \col{panel_sets} and
\col{sources}, not in the evidence columns. Null means ``not measured by that source'', never false.
Every comparison in Section~\ref{sec:compared} is over single-token rows and states its population.

\subsection{The corpora}

Two corpora record containment on single-token rows: the Google Books English 1-grams of the 2020
release (\col{gb_type}) and the English list of wordfreq~3.1.1 (\col{wordfreq_scored}). A feature
layer records how often a string occurs in Google Books. The CEOL's re-stream of Google Books records
containment after admission and admits no string. wordfreq does admit strings:
the census read wordfreq's English list as one of its sources, and the audit consulted it
(Section~\ref{sec:construction}). \cv{wf_census_rows} rows carry the intake channel
\col{census:wordfreq} and \cv{wf_audit_rows} the channel \col{audit:wordfreq}, and wordfreq is the
only source named in the \col{sources} column of \cv{wf_sole_rows} rows.

\subsection{Panel families and evidence tiers}

The \cv{published_set_count} published lists are grouped in \cv{family_count} \emph{panel families}:
WordNet; Wiktionary's English pages; hunspell, whose \cv{hunspell_dictionaries} English dictionaries
(American, British, Canadian, and Australian) count as one family; ENABLE; and SCOWL. A family lists a
row when any of its lists does, and it cannot hold a row when none of its lists can.

The panel families are not independent sources. The American, Canadian, and Australian hunspell
dictionaries are made from SCOWL, as their READMEs state~\cite{libreofficecontributors2026}.
SCOWL's notice names WordNet as an input to its inflection data and names ENABLE's lists among its
own~\cite{atkinson2020}, and ENABLE describes itself as reflecting changes to the Official Scrabble
Players Dictionary~\cite{cooper2026}. Counting SCOWL and hunspell as one family leaves
\cv{strict_core} single-token rows (\cv{strict_core_pct}) listed by two or more families,
\cv{strict_core_less} fewer than the count below; counting ENABLE with them as well leaves
\cv{strict2_core}. Every figure that counts families uses the
\cv{family_count}-family grouping unless it says otherwise.

Each single-token row falls in one of \cv{tier_count} \emph{evidence tiers}
(Table~\ref{tab:tiers}): listed by two or more families; by one; by none, with a corpus holding it
(\emph{corpus only}); or by none and no corpus (\emph{no evidence}). A family that cannot hold a row
does not list it, so the tiers count listings only. A corpus holds a row when it is a
Google Books type or wordfreq scores it: of the corpus-only rows, \cv{co_both} have both signals,
\cv{co_gb_only} Google Books alone, and \cv{co_wf_only} wordfreq alone. The tiers are reported
separately, not averaged, because the corpus-only tier holds \cv{tier_corpus_pct} of the single-token
rows.

\begin{table}[t]
\centering
\caption{The evidence tiers of the \cv{rows_single} single-token rows.}
\label{tab:tiers}
\begin{tabular}{lrr}
\toprule
Evidence tier & Rows & Share of single-token rows (\%) \\
\midrule
2+ families & 417,864 & 4.2 \\
1 family & 672,104 & 6.7 \\
corpus only & 8,032,981 & 80.6 \\
no evidence & 846,459 & 8.5 \\
\bottomrule
\end{tabular}

\end{table}

\subsection{The identifier and the canary}

A row's \col{id} is the first 16 hexadecimal digits of the sha256 of the string's UTF-8 bytes, a
64-bit prefix; every \col{id} in the release is distinct. A name-based UUID is the same for the same
name within a fixed namespace~\cite{rfc9562}; the \col{id} depends on the exact normalized string
alone, so a reader can join
versions by \col{id} or by \col{word}, the data dictionary's join key. The \col{sample} configuration, the rows whose \col{id} begins with
\cv{sample_prefix}, holds the same strings in every version that keeps them.

The release carries a marker for a later test: the dataset card carries a canary string, and the
release record and every Parquet file's metadata carry its GUID. A test of whether a model recalls the
marker can show exposure to some copy of the documentation or files that carried it; it cannot tell
which copy or version, and the release record dates the version separately.

\section{Sources, licenses, and the public projection}
\label{sec:projection}

The build measured \cv{panel_set_count} panel sets. Each source's terms were read at the source, and
the author ruled on each reading. The project withholds the membership of \cv{withheld_list_count}
lists from the public version, for these reasons: NASPA Word List 2023 and Collins Scrabble Words 2021,
which are proprietary; the legacy SOWPODS list, for which no license was found; and Webster's Second
International word lists (web2 and its list of hyphenated terms and phrases, web2a), taken from
the 1934 dictionary, two of whose 1934 copyright registrations were renewed in
1961~\cite{libraryof1961}. The project withholds them; the license review, not this paper, decides
what their terms allow. Their columns and labels are absent. The withholding has these measured
effects:

\begin{itemize}
\item \textbf{Restricted rows.} \cv{rows_restricted} of the pin's \cv{rows_pin} candidates are listed
  only by the two proprietary lists. The private pin holds them; every version built from it leaves
  them out.
\item \textbf{Held-back rows.} Of the \cv{rows_shippable} shippable rows, \cv{held_back_rows} are held
  back in the private record (Table~\ref{tab:heldback}). \cv{hb_tier2} have no published evidence
  once the withheld lists and channels are taken away, and no corpus holds any of them. The other
  \cv{hb_ws} are rows of the census or the audit whose every source label names a withheld list.
  Every one of them has published evidence: \cv{hb_ws_gb} are Google Books types, and
  \cv{hb_ws_wordnet} are listed by WordNet, which is not one of the census's admission sources (\cv{hb_ws_wordnet_gb} are
  both). They are held back because their empty \col{sources} would show that a withheld list holds
  them.
\item \textbf{Recomputed fields.} \col{panel_sets} and \col{panel_count} count the
  \cv{published_set_count} published lists alone, so a row with \col{panel_count} 0 may still be listed
  by a withheld list. \col{provenance_class} is recomputed over the retained channels, the published
  lists, and both corpus signals.
\item \textbf{Withheld feature columns.} \cv{withheld_feature_count} of the feature layer's columns are
  withheld. \cv{withheld_reference_feature_count} were computed against a reference that holds
  Webster's Second, so a value could reveal whether that list holds a string's neighbor, stem, or
  segment. The other \cv{withheld_model_feature_count} come from a character model fitted on a pool
  of strings that at least two-thirds of their applicable panel sets list, the withheld lists among
  them. A test run on version 2.1 fitted the same model on the published lists alone and took the
  difference of the two scores. Among single-token rows that \cv{leak_applicable} of the panel sets can
  hold and \cv{leak_strata_min} to \cv{leak_strata_max} published lists list, the difference ranks
  membership in that pool, whose training partition was used for the fit, above chance: per-character
  AUCs from \cv{leak_auc_min} to \cv{leak_auc_max} over version 2.1's shipped single-token rows in
  those strata, and from \cv{leak_auc_ho_min} to \cv{leak_auc_ho_max} when the positives are
  restricted to the pool's test partition, with ties counted half. Pool membership depends on
  withheld listings, so both columns are withheld.
\end{itemize}

\begin{table}[t]
\centering
\caption{Why a shippable row is held back.}
\label{tab:heldback}
\begin{tabular}{lr}
\toprule
Why the row is held back & Rows \\
\midrule
every listing a withheld list, and no public evidence left & 27,458 \\
a withheld channel label its only evidence & 0 \\
a census or audit row whose every source is a withheld list & 2,558 \\
\bottomrule
\end{tabular}

\end{table}

The build's checks constrain direct disclosure on the shipped rows. No file carries a withheld column
or a held-back row, no file with label columns carries a withheld label, and a guard refuses the
release if one withheld list lists \cv{reach_threshold} or more of the shipped rows of any published
channel label, or of any layer's shipped rows whose sources come out empty. The checks do not prove
that no published field can be combined to infer a withheld listing.

The universe is licensed CC~BY-SA~4.0, the share-alike floor of its Wiktionary-derived fields, and its
\col{NOTICES.md} carries the notices and credits the project identified for its sources; the license
review checks them.

\section{Construction}
\label{sec:construction}

Each step has a committed record, and each count is taken at the revision of the English OpenList
the step names.

\begin{enumerate}
\item \textbf{The English OpenList's universe.} The English OpenList (EOL) is a word list with its own
  grammar and governance. At its revision \cv{archival_eol_revision}, its valid and invalid lists held
  \cv{arch_eol_universe} strings between them.
\item \textbf{The census.} The pinned sources, read again under the grammar, gave
  \cv{census_candidates} candidates in the exception classes that the EOL's letters-only rule could
  not hold: hyphens, accents, a diaeresis, apostrophes, and ligatures.
\item \textbf{The merge.} The two sets merged gave \cv{arch_merge_total} rows, one per string. The
  merge's inherited rows, split by their EOL flag, reproduce the EOL's valid and invalid lists at that
  revision byte for byte.
\item \textbf{The audit.} The audited sources' coverage of the merge was reconciled against the EOL's
  research frame of \cv{audit_frame_size} letters-only strings. Of \cv{audit_misses} strings a source
  listed and the merge lacked, \cv{c4_additions} were added (\cv{audit_in_frame_added} within the
  frame, and \cv{audit_beyond_frame_added} SCOWL forms the EOL's ingest had not read), and
  \cv{audit_wordfreq_left_out} that only wordfreq listed, outside the frame, were left out.
\item \textbf{The re-pin.} At the EOL's later revision \cv{eol_revision}, its lists hold
  \cv{eol_universe_rows} strings: the \cv{repin_left} inherited strings in neither list at that
  revision left the pin. The result is the single-token pin, \cv{v1_rows} candidates.
\item \textbf{The spaced admission.} \cv{pin_spaced} multiword candidates from WordNet and Wiktionary
  were admitted under the grammar's space rule. The pin has \cv{rows_pin} rows, and its single-token
  projection matches the single-token pin's canonical row serialization.
\item \textbf{The layers.} No row is added. Each single-token row is measured against the panel lists
  (the evidence matrix) and given its corpus features (the feature layer); each multiword row carries
  its WordNet and Wiktionary listing; and every row is given its record of how it came in (the
  intake-provenance layer).
\item \textbf{The model-and-engine column:} \col{engine_generated} where the intake kind is
  generated, and \col{none} elsewhere. The EOL's valid dictionary at the pin
  (\cv{model_basis_entries} entries) names a model on none of its entries.
\item \textbf{The release.} The restricted rows are dropped and the public projection applied
  (Section~\ref{sec:projection}), leaving \cv{rows_all} rows in \cv{config_count} configurations
  (Table~\ref{tab:configs}). On every file, the build checked that none of the
  \cv{restricted_words_checked} restricted strings and no withheld column or held-back row is present,
  and, on the files with label columns, no withheld label. On the \col{universe} and \col{features}
  files, it checked that every \col{id} is distinct and that none of \cv{controls_checked} private
  synthetic controls is present. On \col{universe}, it checked that \col{provenance_class} re-derives
  from the file's own fields and that \col{model_or_engine} agrees with its projected rule, which reads
  \col{intake_kind}, \col{eol}, and the private model-basis layer.
  It checked that each of the \cv{filter_count} filter configurations is exactly its predicate's filter
  of \col{universe}.
\end{enumerate}

The release record, \col{pipeline/research/release/release_v1.0.json}, holds every file's rows,
columns, bytes, and sha256, the recorded results of the build's checks, and the reach guard's
threshold; the reach counts themselves are in a private receipt. It names the build's commit,
\cv{code_commit}, and records that tracked files had changed in the working tree when the build ran,
so that commit alone does not identify the code that built the files.

\begin{table}[t]
\centering
\caption{The \cv{config_count} configurations of version~\cv{version}.}
\label{tab:configs}
\small
\begin{tabular}{lrl}
\toprule
Configuration & Rows & What it holds \\
\midrule
\col{universe} & \cv{rows_universe} & every shipped row \\
\col{features} & \cv{rows_features} & the feature layer of the single-token rows \\
\col{panel-listed} & \cv{rows_panel_listed} & rows a published list lists \\
\col{eol-valid} & \cv{rows_eol_valid} & rows of the EOL's valid list at the pin \\
\col{spaced} & \cv{rows_spaced} & multiword rows \\
\col{exceptions} & \cv{rows_exceptions} & rows with any of the census's exception flags \\
\col{sample} & \cv{rows_sample} & rows whose \col{id} begins with \cv{sample_prefix} \\
\bottomrule
\end{tabular}
\end{table}

\section{Contents}
\label{sec:contents}

\subsection{Form}

\paragraph{Length.} The median single-token row has \cv{len_corpus_p50} characters in the corpus-only
tier, \cv{len_two_plus_p50} in the tier listed by two or more families, \cv{len_one_p50} in the tier
listed by one, and \cv{len_none_p50} in the tier with no evidence (Table~\ref{tab:lengths}).
Multiword rows have a median of \cv{len_multiword_p50} characters.

\begin{table}[t]
\centering
\caption{Lengths in characters, by evidence tier.}
\label{tab:lengths}
\small
\begin{tabular}{lrrrrr}
\toprule
Group & Rows & 10th percentile & Median & 90th percentile & Longest \\
\midrule
2+ families & 417,864 & 6 & 9 & 13 & 45 \\
1 family & 672,104 & 7 & 10 & 15 & 100 \\
corpus only & 8,032,981 & 6 & 8 & 12 & 58 \\
no evidence & 846,459 & 7 & 10 & 15 & 63 \\
multiword & 294,865 & 9 & 14 & 21 & 79 \\
\bottomrule
\end{tabular}

\end{table}

\paragraph{Characters.} All \cv{chars_total} characters of the charset occur. Of the accented letters
and ligatures, the least used occurs in \cv{char_rows_nonbase_min} rows and the most used in
\cv{char_rows_nonbase_max}. The apostrophe occurs in \cv{char_rows_apostrophe} rows, the hyphen in
\cv{char_rows_hyphen}, and the space in the \cv{char_rows_space} multiword rows.

\paragraph{Classes.} \cv{flag_plain} rows (\cv{flag_plain_pct}) are plain. The other \cv{flag_any}
carry at least one class, in \cv{class_combos_nonplain} combinations (Table~\ref{tab:flags}).

\begin{table}[t]
\centering
\caption{Rows by class flag. A row can carry several flags.}
\label{tab:flags}
\begin{tabular}{lrr}
\toprule
Class flag & Rows & Share of all rows (\%) \\
\midrule
\texttt{accented} & 5,245 & 0.05 \\
\texttt{apostrophe-form} & 140,851 & 1.37 \\
\texttt{contraction} & 1,841 & 0.02 \\
\texttt{diaeresis} & 1,530 & 0.01 \\
\texttt{hyphenated} & 63,421 & 0.62 \\
\texttt{ligature} & 2,002 & 0.02 \\
\texttt{s-final} & 137,018 & 1.33 \\
\texttt{spaced} & 294,865 & 2.87 \\
\bottomrule
\end{tabular}

\end{table}

\paragraph{Apostrophes.} Of the \cv{apos_rows} rows with an apostrophe, \cv{flag_s_final} have a token
ending in an apostrophe and s (in \cv{s_final_token_only} of them, only tokens before the last),
\cv{flag_contraction} carry the contraction flag, \cv{apos_both_flags} carry both flags, and
\cv{apos_neither} carry neither.

\paragraph{Multiword terms.} WordNet and Wiktionary both list \cv{mw_both} of the multiword rows,
WordNet alone \cv{mw_wordnet_only}, and Wiktionary alone \cv{mw_wiktionary_only}. Every multiword row
is listed by at least one of them, because they are its only sources.

\paragraph{Variants.} \cv{var_groups} groups of two or more rows (\cv{var_rows} rows) are equal once
hyphens, spaces, apostrophes, and combining accent marks are removed. The largest group has \cv{variant_largest}
members. Table~\ref{tab:variants} shows the \cv{variant_patterns_shown} most common of
\cv{variant_patterns} patterns of which kinds of mark are present in some members and absent in
others; the pattern does not compare a mark's count or position. Of the \cv{pairs_total} groups of
two, the two members share a published list in \cv{pairs_share}, are both listed but by no common
list in \cv{pairs_both_no_common}, have one member listed in \cv{pairs_one}, and have neither listed
in \cv{pairs_none}.

\begin{table}[t]
\centering
\caption{Variant groups by which kinds of mark are present in some members and absent in others;
the \cv{variant_patterns_shown} most common of \cv{variant_patterns} patterns.}
\label{tab:variants}
\begin{tabular}{lr}
\toprule
Varies by & Groups \\
\midrule
apostrophe & 87,070 \\
space & 55,617 \\
hyphen & 24,596 \\
hyphen, space & 14,078 \\
accent & 4,191 \\
hyphen, apostrophe & 2,000 \\
space, apostrophe & 1,766 \\
hyphen, space, apostrophe & 1,218 \\
apostrophe, accent & 785 \\
no kind of mark changes presence & 309 \\
\bottomrule
\end{tabular}

\end{table}

\subsection{Lexical category}

The universe records no part of speech, lemma, or inflection. WordNet lists \cv{list_wordnet_listed}
single-token rows and \cv{mw_wordnet} multiword rows, and those rows could take a part of speech from
WordNet's synsets.

\subsection{Frequency}

wordfreq scores few rows: \cv{zipf_with_two_plus_pct} of the tier listed by two or more families,
\cv{zipf_with_one_pct} of the tier listed by one, and \cv{zipf_with_corpus_pct} of the corpus-only
tier. Google Books holds more: \cv{gb_with_two_plus_pct}, \cv{gb_with_one_pct}, and
\cv{gb_with_corpus_pct} of the same tiers. Over the rows that have a value, the median Zipf value is
\cv{zipf_p50_two_plus} in the first tier, \cv{zipf_p50_one} in the second, and
\cv{zipf_p50_corpus} in the corpus-only tier, and the median $\log_{10}$ count of Google Books matches
is \cv{gb_p50_two_plus}, \cv{gb_p50_one}, and \cv{gb_p50_corpus}.

By the number of families that list a row, the median $\log_{10}$ Google Books count, among rows
with a Google Books value, rises from \cv{u5_gb_p50_k0} for rows no family lists to \cv{u5_gb_p50_k5}
for rows all \cv{family_count} families list, a group of \cv{u5_rows_k5} rows (Table~\ref{tab:u5}). The span of years over which a string occurs in
Google Books has a median, over the Google Books types, of \cv{span_p50_two_plus} years in the first
tier, \cv{span_p50_one} in the second, and \cv{span_p50_corpus} in the corpus-only tier. It is a span in
a corpus of scanned books, not a date of first use, and the corpus counts occurrences without
weighting them by readership~\cite{pechenick2015}.

\begin{table}[t]
\centering
\caption{Single-token rows by the number of families that list them. The row counts cover every
row; each median covers the rows that have that value.}
\label{tab:u5}
\small
\begin{tabular}{rrrr}
\toprule
Families listing & Rows & Median Zipf & Median log10 Google Books matches \\
\midrule
0 & 8,879,440 & 1.35 & 2.51 \\
1 & 672,104 & 1.60 & 2.96 \\
2 & 217,093 & 1.62 & 3.21 \\
3 & 76,918 & 1.66 & 3.84 \\
4 & 75,445 & 1.97 & 4.78 \\
5 & 48,408 & 2.63 & 5.89 \\
\bottomrule
\end{tabular}

\end{table}

\subsection{Provenance}

Every row records its \emph{intake layer} (the step of Section~\ref{sec:construction} that brought it
in), its \emph{intake kind} (obtained from a source, derived from pinned sources, or generated), and
its \emph{intake channels}. A channel records how a string entered the pool and is never an
attestation. For a string inherited from the EOL, the channel is the intake its EOL entry publishes
or, where none is published, a fallback: the EOL list it came in with or the EOL's candidate pool,
which does not name a particular upstream source; for the CEOL's own layers, it
names the source in the census, audit, spaced, or ledger namespace. A row can carry several channels, so the
channel counts overlap (Table~\ref{tab:channels}).

The EOL's invalid list contributes \cv{p1_eol_invalid} rows, its valid list \cv{p1_eol_valid}, the
census \cv{p1_c1_exceptions}, the audit \cv{p1_c4_additions}, and the spaced admission
\cv{p1_spaced_c015}. The census's other \cv{census_not_shipped} candidates and the audit's other
\cv{audit_not_shipped} additions are restricted or held back (Section~\ref{sec:projection}). The largest channel, \col{eol:candidate-pool}, is a channel of
\cv{p4_top_rows} rows, of which a published family lists \cv{p4_top_listed_pct}.

\begin{table}[t]
\centering
\caption{The \cv{channels_shown} largest of \cv{p4_channels} intake channels. A row can carry several. The
corpora do not measure multiword rows.}
\label{tab:channels}
\small
\begin{tabular}{lrrr}
\toprule
Intake channel & Rows & Listed by a family (\%) & In a corpus (\%) \\
\midrule
\texttt{eol:candidate-pool} & 9,009,717 & 5.7 & 92.7 \\
\texttt{eol:synthetic-engine} & 330,363 & 1.2 & 2.6 \\
\texttt{eol:list} & 312,944 & 84.8 & 94.1 \\
\texttt{spaced:wiktionary\_english} & 250,893 & 100 & not measured \\
\texttt{audit:wiktionary\_english} & 104,848 & 100 & 80.0 \\
\texttt{audit:scowl} & 90,637 & 100 & 80.8 \\
\texttt{census:hunspell\_en\_GB} & 86,480 & 100 & 16.2 \\
\texttt{census:scowl} & 79,418 & 100 & 10.3 \\
\texttt{spaced:wordnet} & 63,483 & 100 & not measured \\
\texttt{census:wiktionary\_english} & 56,104 & 100 & 8.3 \\
\texttt{census:hunspell\_en\_AU} & 33,293 & 100 & 32.2 \\
\texttt{census:hunspell\_en\_US} & 33,276 & 100 & 32.3 \\
\bottomrule
\end{tabular}

\end{table}

\paragraph{The synthetic engine.} \cv{generated_rows} rows (\cv{generated_pct}) were made by the EOL's
synthetic engine. Its code does not survive, and its own per-row record of the method is empty, but
its author describes the methods: affixes applied without a part-of-speech filter to the list as it
then stood; chemical names varied by substituent, and short chains of root, infix, and suffix, made
by substitution over a hand-written list; medical and scientific prefixes; and Greek plurals. Each
row carries the method that a heuristic infers from its form where it can (\cv{p2_none} rows are
unresolved; Appendix~\ref{app:tables}). The methods are the author's account and were not checked
against the engine itself. A published family lists
\cv{p2_listed} of these rows (\cv{p2_listed_pct}).

\paragraph{Provenance classes.} Over the published fields, \cv{prov_channel_and_containment} rows have
both a recorded channel and containment in a list or a corpus, \cv{prov_channel_only} a channel only,
\cv{prov_containment_only} containment only, and \cv{prov_neither} neither.

\section{The sources compared}
\label{sec:compared}

This section covers the single-token rows.

\paragraph{Methods.} The main pair measure is directed containment: of the strings list A lists that
list B could hold, the share that B lists. It is computed over the rows both lists can hold. Agreement
is counted as $k$ of $m$: the families that list a row out of the families that could. The analysis
is descriptive: it reports exact counts of the release's rows and runs no statistical tests. It
makes no estimate of how many English strings no list holds. The descriptive analysis
(\col{pipeline/research/analysis/descriptive.py}) defines the populations of its counts.

\subsection{Sizes and blind classes}

Table~\ref{tab:lists} gives each list's size as a share of all single-token rows, not of the rows it
could hold. Wiktionary lists the most single-token rows, \cv{list_wiktionary_listed}, and is blind to
no class. WordNet lists the fewest, \cv{list_wordnet_listed}. ENABLE cannot hold \cv{list_enable_blind}
rows, every row with an exception class, and each hunspell dictionary but the British cannot hold
\cv{list_hunspell_us_blind}.

\begin{table}[t]
\centering
\caption{The published lists over the \cv{rows_single} single-token rows.}
\label{tab:lists}
\small
\begin{tabular}{lrrr}
\toprule
List & Listed & Cannot hold & Listed (\% of single-token rows) \\
\midrule
WordNet & 80,441 & 8,147 & 0.81 \\
Wiktionary (English) & 831,973 & 0 & 8.35 \\
hunspell en\_US & 121,252 & 61,747 & 1.22 \\
hunspell en\_GB & 262,413 & 1,940 & 2.63 \\
hunspell en\_CA & 120,947 & 61,747 & 1.21 \\
hunspell en\_AU & 121,490 & 61,747 & 1.22 \\
ENABLE & 172,823 & 187,266 & 1.73 \\
SCOWL & 518,438 & 55,759 & 5.20 \\
\bottomrule
\end{tabular}

\end{table}

\subsection{Containment}

Table~\ref{tab:containment} gives, for each pair, the share of A's rows that B can hold which B lists;
each share below is among the rows the receiving list can hold. The asymmetry shows nesting. Of ENABLE's rows that SCOWL can hold, SCOWL lists \cv{enable_in_scowl},
and of SCOWL's rows that ENABLE can hold, ENABLE lists \cv{cont_scowl_in_enable}: ENABLE is nearly a
subset of SCOWL. Of the American hunspell dictionary's rows, the Canadian lists
\cv{cont_hunspell_us_in_hunspell_ca}. Of WordNet's rows, SCOWL lists \cv{cont_wordnet_in_scowl} and
Wiktionary \cv{cont_wordnet_in_wiktionary}, while of Wiktionary's rows, WordNet lists
\cv{cont_wiktionary_in_wordnet} and SCOWL \cv{cont_wiktionary_in_scowl}: Wiktionary holds much that
the others do not.

\begin{table}[t]
\centering
\caption{Directed containment (\%): of the rows list A lists that list B can hold, the share B lists.
Shares are rounded to whole percent, except where that would print 0 or 100 for a share that is
neither.}
\label{tab:containment}
\small
\begin{tabular}{lrrrrrrrr}
\toprule
A $\downarrow$, B $\rightarrow$ & WN & Wikt & US & GB & CA & AU & EN & SC \\
\midrule
WN & --- & 84 & 57 & 69 & 57 & 57 & 71 & 94 \\
Wikt & 8 & --- & 10 & 18 & 10 & 10 & 21 & 43 \\
US & 36 & 66 & --- & 88 & 99 & 97 & 86 & 88 \\
GB & 21 & 56 & 42 & --- & 42 & 43 & 62 & 69 \\
CA & 37 & 66 & 99 & 89 & --- & 98 & 86 & 89 \\
AU & 36 & 66 & 97 & 90 & 98 & --- & 85 & 89 \\
EN & 32 & 96 & 44 & 64 & 44 & 43 & --- & 99.7 \\
SC & 14 & 65 & 21 & 34 & 21 & 21 & 39 & --- \\
\bottomrule
\end{tabular}

\end{table}

\subsection{Combinations and what each family holds alone}

Over the \cv{s3_rows} rows every family can hold, the rows fall in \cv{combos_total} combinations of
families. \cv{m5_k0} (\cv{m5_k0_pct}) are listed by no family. Table~\ref{tab:s3} shows the
\cv{combos_shown} largest combinations. Over all single-token rows, Table~\ref{tab:alone} gives the
rows that only one family lists: Wiktionary alone lists \cv{alone_wiktionary_pct} of the
\cv{listed_single} listed single-token rows, and ENABLE alone lists \cv{alone_enable}. These counts
include rows that another family cannot hold; over the rows all \cv{family_count} families can hold,
Wiktionary alone lists \cv{wikt_alone_m5} of the \cv{m5_listed} listed rows
(\cv{wikt_alone_m5_pct}).

\begin{table}[t]
\centering
\caption{The \cv{combos_shown} largest of \cv{combos_total} combinations, over the rows every family can
hold.}
\label{tab:s3}
\small
\begin{tabular}{lr}
\toprule
Families that list the row & Rows \\
\midrule
none & 8,872,182 \\
Wiktionary & 426,798 \\
Wiktionary, SCOWL & 141,928 \\
SCOWL & 84,810 \\
Wiktionary, hunspell, ENABLE, SCOWL & 65,712 \\
WordNet, Wiktionary, hunspell, ENABLE, SCOWL & 48,408 \\
Wiktionary, ENABLE, SCOWL & 46,033 \\
hunspell & 27,682 \\
Wiktionary, hunspell, SCOWL & 18,713 \\
Wiktionary, hunspell & 11,824 \\
WordNet, SCOWL & 6,381 \\
WordNet, Wiktionary, ENABLE, SCOWL & 6,076 \\
\bottomrule
\end{tabular}

\end{table}

\begin{table}[t]
\centering
\caption{Rows only one family lists, over all single-token rows.}
\label{tab:alone}
\small
\begin{tabular}{lrr}
\toprule
Family & Rows only it lists & Share of listed single-token rows (\%) \\
\midrule
WordNet & 1,522 & 0.1 \\
Wiktionary & 476,575 & 43.7 \\
hunspell & 72,784 & 6.7 \\
ENABLE & 78 & <0.1 \\
SCOWL & 121,145 & 11.1 \\
\bottomrule
\end{tabular}

\end{table}

\subsection{Agreement}

Among the rows all \cv{family_count} families can hold, \cv{m5_listed} are listed by at least one. Of
those, \cv{m5_k1_pct_of_listed} are listed by exactly one family and \cv{m5_k5_pct_of_listed} by all
\cv{family_count} (Table~\ref{tab:kofm}).

\begin{table}[t]
\centering
\caption{Single-token rows by the number of families that can hold them ($m$) and the number that
list them ($k$).}
\label{tab:kofm}
\small
\begin{tabular}{rrrrrrr}
\toprule
$m$ & $k=0$ & $k=1$ & $k=2$ & $k=3$ & $k=4$ & $k=5$ \\
\midrule
1 & 14 & 1,926 &  &  &  &  \\
2 &  & 212 & 7 &  &  &  \\
3 & 2,418 & 50,968 & 4,810 & 1,392 &  &  \\
4 & 4,826 & 78,113 & 42,376 & 190 & 14 &  \\
5 & 8,872,182 & 540,885 & 169,900 & 75,336 & 75,431 & 48,408 \\
\bottomrule
\end{tabular}

\end{table}

\subsection{Coverage by frequency}

Table~\ref{tab:coverage} gives, for each band of wordfreq's Zipf value, the share of the rows a list
could hold that it lists. Rows that wordfreq does not score form their own band, which holds most
rows. By this measure, coverage rises from band to band across the scored bands for every list, in
line with what Michel and colleagues found for two dictionaries over bands of the Google Books
lexicon~\cite{michel2011}. In the band of Zipf 5 and over, the American hunspell dictionary lists
\cv{cov_hunspell_us_b5plus} and WordNet \cv{cov_wordnet_b5plus}. Weighted by wordfreq's frequency
estimate (each row's weight over the weights of the band's rows the list can hold), WordNet's share in
that band falls to \cv{covw_wordnet_b5plus} (Appendix~\ref{app:tables}). One possible reason, not
measured here, is that WordNet holds only nouns, verbs, adjectives, and adverbs, while frequent strings
include function words.

\begin{table}[t]
\centering
\caption{Coverage by Zipf band (\%), by strings: of the rows in a band that a list can hold, the share
it lists. Rounded as in Table~\ref{tab:containment}.}
\label{tab:coverage}
\small
\begin{tabular}{lrrrrrr}
\toprule
List & no value & 1 to 2 & 2 to 3 & 3 to 4 & 4 to 5 & 5 and over \\
\midrule
WordNet & 0.2 & 9 & 31 & 59 & 75 & 83 \\
Wiktionary (English) & 7 & 32 & 57 & 82 & 94 & 98 \\
hunspell en\_US & 0.4 & 14 & 51 & 90 & 99 & 99.9 \\
hunspell en\_GB & 1 & 25 & 64 & 93 & 99 & 99.8 \\
hunspell en\_CA & 0.4 & 14 & 51 & 89 & 99 & 99.8 \\
hunspell en\_AU & 0.4 & 14 & 51 & 89 & 99 & 99.7 \\
ENABLE & 1 & 20 & 49 & 77 & 90 & 94 \\
SCOWL & 4 & 29 & 59 & 86 & 97 & 99 \\
\bottomrule
\end{tabular}

\end{table}

\subsection{Varieties, the core, and the tails}

\paragraph{The American and British dictionaries.} Over the single-token rows both dictionaries can
hold,
\cv{us_not_gb} rows are in the American hunspell dictionary and not the British, \cv{gb_not_us} in the
British and not the American, and \cv{us_and_gb} in both. Another \cv{gb_listed_us_blind} rows are
in the British dictionary and of a class the American cannot hold. Of the
\cv{hun_all4_hold_listed} rows that all \cv{hunspell_dictionaries} dictionaries can hold and at least
one lists, \cv{hun_all4_hold_some_pct} are listed by some and not all; over every row at least one of
them lists, the share is \cv{hun_any_some_pct}.

\paragraph{The consensus core.} \cv{core_rows} rows (\cv{core_pct}) are listed by two or more
families. The EOL's valid list contributes \cv{core_layer_eol_valid} of them and its invalid list
\cv{core_layer_eol_invalid}: strings on the EOL's invalid list that two or more panel families
list.

\paragraph{The tails.} \cv{freq3_unlisted} rows with a Zipf value of 3 or more are listed by no
family. In the other direction, \cv{listed_no_corpus} listed rows (\cv{listed_no_corpus_pct} of the
listed single-token rows) are in neither corpus.

\subsection{What each list adds}

As lists are added in every order, the union of their listed rows grows from a mean of
\cv{acc_mean_1} rows for one list (rounded) to \cv{acc_mean_8} for all \cv{published_set_count}; at the
intermediate steps, the order matters, as the range shows (Appendix~\ref{app:tables}).

\section{Limitations}
\label{sec:limitations}

\begin{itemize}
\item \textbf{Most rows rest on a single corpus.} \cv{inh_unlisted_gb} single-token rows
  (\cv{inh_unlisted_gb_pct_single} of the single-token rows; \cv{inh_unlisted_gb_pct_all} of all rows)
  were inherited from the EOL, are listed by none of the \cv{published_set_count} published lists, and
  are Google Books types. For \cv{inh_unlisted_gb_only} of them, wordfreq does not score the string, so
  Google Books is their only published containment, and it is a corpus built from scanned text, which
  carries recognition errors. In all, \cv{limit_gb_all} rows are Google Books types.
\item \textbf{\cv{limit_no_panel_no_corpus} single-token rows are listed by no published list and are
  in neither corpus.} What supports them is their intake channel, which records how the string entered
  the pool and attests nothing.
\item \textbf{\cv{generated_rows} rows (\cv{generated_pct}) were made by the EOL's synthetic engine}
  (Section~\ref{sec:contents}). \col{model_or_engine} marks them \col{engine_generated}.
\item \textbf{Some rows of the EOL's valid list were decided with a language model,} by the EOL's
  verification pipeline or by a model's verdict alone. No row says which.
\item \textbf{The withholding changes what the panel columns count} (Section~\ref{sec:projection}).
\item \textbf{Four Google Books features sum over raw forms.} The Google Books
  aggregate sums, over the raw forms that normalize to one string (case variants among them), the
  match counts, the volume counts, and the number of years each form occurs in, and it caps only that
  last sum at the span. A year two forms share counts twice before the cap, so \col{gb_span_fill} can
  exceed the share of distinct years in the span, and \col{gb_volumes_per_year} divides the summed
  volume count by the capped year count. \col{gb_log10_volume} and \col{gb_tokens_per_volume} use the
  summed volume count. The data dictionary states this; a
  later version can count distinct years.
\item \textbf{\col{wordfreq_scored} is false, not null, on the \cv{hyph_wf_false} hyphenated
  single-token rows:} the CEOL's feature code does not score hyphenated strings, so
  \col{wordfreq_zipf} is NaN on them.
\item \textbf{The variety and properness columns are proxies.} \col{variety_gb_not_us} and
  \col{variety_us_not_gb} compare two hunspell dictionaries, and \col{wn_instance},
  \col{wn_capitalized}, and \col{wk_proper_noun} record what WordNet and Wiktionary say about a
  multiword row; none is a variety label or a proper-name filter.
\item \textbf{The \col{sources} column has two meanings:} a per-source list on the CEOL's own rows,
  and the dataset-level tag \col{eol_universe@c550751e} on inherited rows, whose per-source listing is
  in the evidence columns.
\item \textbf{The Croissant file} declares the Responsible AI namespace but carries none of its
  properties, and it has no publication date.
\item \textbf{\col{eol} is a fact at one revision of the EOL,} \cv{eol_revision}; the EOL changes as
  it takes words in.
\end{itemize}

\paragraph{What the universe does not record:} part of speech; lemma and inflection relations;
definitions and senses; pronunciation, syllables, etymology, and a first attestation year; variety and
register labels; a frequency drawn from many kinds of text for most rows; and validity, by design.

\section{Uses and out-of-scope uses}
\label{sec:uses}

Among the resources in Section~\ref{sec:related} and Appendix~\ref{app:resources}, none stores a
per-string listing for each of several pinned word lists, with ``cannot hold'' marked apart from
``not listed'' on single-token rows. KoAGE's evidence packets for its verification stage carry the temporal exclusion status of a candidate and its parts~\cite{kim2026koage}. The
universe supports these uses, with their caveats:

\begin{enumerate}
\item \textbf{A ``listed by 0 of $N$ lists'' screen for invented stimuli,} never a ``not a word''
  verdict. For a single-token row, $N$ is the number of published lists that can hold it, at most
  \cv{published_set_count} in this version.
\item \textbf{Auditing pseudoword generators:} finding generated strings that a list holds, or that
  are variants of a listed string.
\item \textbf{Per-list membership and eligibility,} with ``not measured'' kept apart from ``not
  listed''.
\item \textbf{A reproducible frame for choosing items} in rating and norming studies: a stated filter
  over a pinned version, with its code.
\item \textbf{Frequency matching in the very low range,} with missing values explicit.
\item \textbf{A dated exclusion list} for benchmarks of novel words, run again at each version.
\end{enumerate}

\paragraph{Checking a made-up string.} Put it through the grammar and look it up in the
\col{universe} configuration. A string with no row has no per-list result in this version. For a
single-token row, count the lists that can hold it as well as those that list it, remembering that
$-1$ means the list cannot hold it; treat a listing as a listing, not as use; and treat absence from
every list as absence from the lists, not from text. A multiword row is checked against two lists
only, through \col{panel_sets}.

\paragraph{Out of scope:} calling a string a word or not a word; using the universe as a dictionary;
proper names; and any claim about what a model was trained on.

\section{Maintenance and versions}

A released version never changes; a correction comes as a new version with a change log. A row
leaves a public version only under the rules of the project's upkeep policy, and a version may be
withdrawn, never edited. A request to take down a source or a row is acknowledged within 7 days and,
once verified, acted on within 14 days of its verification, and each version's change log records
every withholding, as counts. A checkpoint that measures every source again and builds a new pin and
version is designed and partly built. A later pin of the hunspell dictionaries may take in the
en-wl/wordlist 2026.02.25 release of them, whose notes say it used language models to help filter
over 1,500 new words from a corpus~\cite{atkinson2026release}.

\section{Responsible use and disclosures}

\begin{itemize}
\item \textbf{Language models.} No language model decided any row the CEOL added. The EOL decided
  some rows of its valid list with a language model, Google Gemini 3 Flash Preview by its author's
  account, and the EOL's own report says the earlier phases of its pipeline used Gemini 2.0 Flash and
  Gemini 2.5 Flash. No row records which rows those are, and the EOL's dictionary at the pinned
  revision names a model on none of its entries. \col{model_or_engine} marks the rows the EOL's synthetic engine made, so a list can
  keep them or leave them out.
\item \textbf{Offensive content.} The universe holds strings its sources hold, and some are
  offensive. This version does not measure their share, because the measure needs an offensive-term
  lexicon whose license allows the use, and the project has not chosen one. This paper prints no
  string of the universe.
\item \textbf{Contamination.} The release's date and file hashes identify it, and its canary
  supports a later test of whether a model recalls the marker. Version 1.0 and the internal version
  2.1 share the marker, and the marker can spread without the rows, so such a test cannot say when or
  how a model met it.
\end{itemize}

\section{Availability and reproducibility}
\label{sec:availability}

Version~\cv{version} is built for public release and is private until a license review
confirms the terms under which its sources may be published. An open subset, the rows listed by
WordNet, Wiktionary, the American, Canadian, and Australian hunspell dictionaries, or SCOWL at sizes
10 to 70, is public as \col{ryanjosephkamp/ceol-universe-open}~\cite{ceolopen2026}. A DOI is planned
for the full release. The code and every committed record are in the project's
repository~\cite{ceolrepo2026}.

The descriptive analysis reruns from the release files, which are private until the full release. From the repository's \col{pipeline}
directory, with the files in place,
\begin{quote}\small\ttfamily python3 -m research.analysis.descriptive --version v1.0 --check\end{quote}
reads the release's \col{universe} and \col{features} files, checks each against its recorded sha256,
reads the committed EOL comparison record, renders \col{analysis.json}, and compares the rendering
with the committed file; without \col{--check}, it writes the file. A rebuild of the release
itself (\col{research.release.build_release}) needs private inputs: the five pinned layers, the
restricted-row list, the intake ledger, and the synthetic controls. This paper did not rebuild the
layers or the release. Beyond the release files, the held-back split of
Section~\ref{sec:projection} also reads the private held-back receipt and two internal layers.

\renewcommand{\bibfont}{\small\raggedright}
\bibliographystyle{unsrtnat}
\bibliography{generated/refs_page,refs_added}

\appendix

\section{Further tables}
\label{app:tables}

\begin{table}[htbp]
\centering
\caption{Coverage by Zipf band (\%), weighted by wordfreq's frequency estimate.}
\small
\begin{tabular}{lrrrrr}
\toprule
List & 1 to 2 & 2 to 3 & 3 to 4 & 4 to 5 & 5 and over \\
\midrule
WordNet & 11 & 38 & 62 & 77 & 57 \\
Wiktionary (English) & 36 & 63 & 85 & 95 & 99 \\
hunspell en\_US & 18 & 61 & 93 & 99 & >99.9 \\
hunspell en\_GB & 31 & 72 & 94 & 99 & >99.9 \\
hunspell en\_CA & 18 & 60 & 93 & 99 & >99.9 \\
hunspell en\_AU & 18 & 60 & 92 & 99 & 99.9 \\
ENABLE & 24 & 56 & 80 & 92 & 93 \\
SCOWL & 33 & 66 & 89 & 98 & 99.9 \\
\bottomrule
\end{tabular}

\end{table}

\begin{table}[htbp]
\centering
\caption{The union of listed single-token rows as lists are added in every order.}
\small
\begin{tabular}{rrrr}
\toprule
Lists & Mean union & Smallest & Largest \\
\midrule
1 & 278,722 & 80,441 & 831,973 \\
2 & 451,746 & 122,249 & 1,013,749 \\
3 & 592,049 & 124,889 & 1,084,450 \\
4 & 713,157 & 160,092 & 1,088,351 \\
5 & 820,725 & 240,688 & 1,089,887 \\
6 & 918,139 & 349,973 & 1,089,965 \\
7 & 1,007,439 & 613,393 & 1,089,968 \\
8 & 1,089,968 & 1,089,968 & 1,089,968 \\
\bottomrule
\end{tabular}

\end{table}

\begin{table}[htbp]
\centering
\caption{The EOL's synthetic engine: rows by the method inferred from their form.}
\small
\begin{tabular}{lrr}
\toprule
Engine, as inferred & Rows & Listed by a family \\
\midrule
affixation/plural & 149,139 & 3,270 \\
affixation/verb-form & 117,166 & 439 \\
affixation/superlative & 30,674 & 74 \\
affixation/comparative-or-agent & 29,933 & 120 \\
greek-plural/-ses & 978 & 23 \\
chemical-variation/substituent-prefix & 904 & 0 \\
chemical-variation/root-chain & 691 & 4 \\
medical-scientific-prefixing/domain-prefix & 350 & 12 \\
unresolved/none & 219 & 5 \\
affixation/general-prefix & 149 & 5 \\
greek-plural/-des & 142 & 1 \\
affixation/gerund-plural & 18 & 0 \\
\bottomrule
\end{tabular}

\end{table}

\begin{table}[htbp]
\centering
\caption{Rows by intake layer and kind, and by evidence tier.}
\small
\resizebox{\linewidth}{!}{\begin{tabular}{llrrrrr}
\toprule
Intake layer & Kind & 2+ families & 1 family & Corpus only & No evidence & Multiword \\
\midrule
\texttt{c1\_exceptions} & derived & 48,789 & 131,219 & 7,240 & 18 & 0 \\
\texttt{c4\_additions} & derived & 71,407 & 56,921 & 147 & 0 & 0 \\
\texttt{eol\_invalid} & generated & 1,757 & 845 & 3,522 & 292,996 & 0 \\
\texttt{eol\_invalid} & obtained & 93,488 & 418,337 & 7,973,107 & 524,785 & 0 \\
\texttt{eol\_valid} & generated & 697 & 654 & 1,918 & 27,974 & 0 \\
\texttt{eol\_valid} & obtained & 201,726 & 64,128 & 47,047 & 686 & 0 \\
\texttt{spaced\_c015} & derived & 0 & 0 & 0 & 0 & 294,865 \\
\bottomrule
\end{tabular}
}
\end{table}

\section{Comparable resources}
\label{app:resources}

Sizes and licenses are as stated in the cited sources, read on October 26, 2026, unless an entry or its
reference gives another date; some pages have changed since (kaikki.org, for one, now names a later extraction of the
same dump). A \emph{game-derived} mark records what a list's documentation, or its maintainer's documentation of the
same source lists, says about taking its selection, in whole or in part, from a proprietary game word list;
it is not a finding about any entry.


\subsection*{Lexicons, Wiktionary, and its extractions}

\resource{Princeton WordNet 3.0}
\resfield{What it is}{A lexical database of English that groups nouns, verbs, adjectives, and adverbs into synonym sets (synsets) linked by conceptual-semantic and lexical relations. Princeton states that it is no longer developed and names Open English WordNet and the Global WordNet Association as community-led resources that build on it~\cite{princetonuniversity2026}.}
\resfield{Size, as stated}{155,287 unique strings counted within each part of speech (147,278 distinct overall), 117,659 synsets, and 206,941 word-sense pairs~\cite{princetonuniversity2026a}.}
\resfield{License, as stated}{The WordNet 3.0 license: use, copying, modification, and distribution for any purpose without fee, provided the copyright notice and statements appear on all copies; Princeton's name may not be used in publicity~\cite{princetonuniversity2026b}.}
\resfield{What it records per entry}{For each lemma and part of speech, its synsets, pointer (relation) types, and a count of tagged senses; for each synset, its words, lexicographer file, relations, and a gloss~\cite{princetonuniversity2026c}.}
\resfield{What it does not record}{Parts of speech other than the 4 above~\cite{princetonuniversity2026a}; the database format has no pronunciation or etymology field~\cite{princetonuniversity2026c}.}

\resource{Open English WordNet, 2025 edition}
\resfield{What it is}{A fork of Princeton WordNet developed under an open-source methodology, released in annual editions since 2019; from the 2025 release, proper nouns are in a separate resource, Open English Namenet~\cite{openenglish2026}.}
\resfield{Size, as stated}{2025 edition: 135,969 words and 107,519 synsets; 2025 Plus edition, which adds curated proper nouns: 161,875 words and 120,564 synsets~\cite{openenglish2026}.}
\resfield{License, as stated}{CC BY 4.0 for the team's work, with the Princeton WordNet license applying to the underlying data; attribution to both is required~\cite{openenglish2026}.}
\resfield{What it records per entry}{Lexical entries (lemma and part of speech) and senses, grouped into synsets that carry a part of speech, an Interlingual Index link, relations, and examples~\cite{openenglish2026}; definitions~\cite{openenglish2026a}; pronunciation for nearly 35,000 entries and Wikidata links for over 6,000, as of the 2021 release~\cite{openenglish2026a}.}
\resfield{What it does not record}{Proper nouns in the 2025 edition, which moved them to Open English Namenet~\cite{openenglish2026}.}

\resource{English Wiktionary}
\resfield{What it is}{A collaboratively edited dictionary whose pages hold one section per language; an entry is a single language section~\cite{wiktionarycontributors2026}. The text is an unstructured wiki maintained by volunteers~\cite{ylonen2022}.}
\resfield{Size, as stated}{English: 1,396,346 entries, 719,082 of them lemma entries, generated from the October 1, 2026 dump~\cite{wiktionarycontributors2026}.}
\resfield{License, as stated}{CC BY-SA 4.0 and the GNU Free Documentation License, under the Wikimedia Terms of Use in effect since June 7, 2023~\cite{wikimediafoundation2023}.}
\resfield{What it records per entry}{Wikitext per language section; the next row lists what an extractor recovers from it~\cite{ylonen2022}.}
\resfield{What it does not record}{A fixed schema: it is an unstructured wiki~\cite{ylonen2022}.}

\resource{Wiktextract and kaikki.org}
\resfield{What it is}{An extractor that expands Wiktionary's templates and Lua modules to produce structured JSON~\cite{ylonen2022}; kaikki.org publishes its output~\cite{kaikkiorg2026}.}
\resfield{Size, as stated}{English: 1,396,111 distinct words, extracted October 21, 2026 from the English Wiktionary dump dated October 1, 2026, as the page stated on October 26, 2026~\cite{kaikkiorg2026}.}
\resfield{License, as stated}{The same licenses as Wiktionary, CC BY-SA and the GFDL, with no version stated~\cite{kaikkiorg2026}.}
\resfield{What it records per entry}{Lemmas, inflected forms, translations, etymology, usage examples, pronunciations (with sound-file URLs), lexical and semantic relations, and morphological, syntactic, semantic, topical, and dialectal annotations~\cite{ylonen2022}.}
\resfield{What it does not record}{A fixed release: the data is re-extracted regularly, usually at least once a week~\cite{kaikkiorg2026}.}

\resource{DBnary}
\resfield{What it is}{Lexical data extracted from 27 Wiktionary language editions and published as linked open data in the OntoLex model, one version per Wikimedia dump~\cite{serasset2026,serasset2026a}.}
\resfield{Size, as stated}{Not confirmed: the site was unreachable on October 26, 2026, and its home page states no counts~\cite{serasset2026}.}
\resfield{License, as stated}{Stated 2 ways: CC BY-SA 3.0 on its site~\cite{serasset2026}; CC BY 4.0 in its Zenodo record for version 20260601~\cite{serasset2026a}.}
\resfield{What it records per entry}{Pages, entries, senses, translations, and lexical relations, with a set of disambiguated translations as an added layer~\cite{serasset2026}.}
\resfield{What it does not record}{Not stated in the pages read.}

\resource{MorphyNet}
\resfield{What it is}{A multilingual database of inflectional and derivational morphology drawn from Wiktionary editions~\cite{morphynetproject2026}.}
\resfield{Size, as stated}{English: 396,772 words with 649,594 inflected forms; 67,412 words with 225,131 derived forms~\cite{morphynetproject2026}.}
\resfield{License, as stated}{CC BY-SA 3.0~\cite{morphynetproject2026}.}
\resfield{What it records per entry}{Inflection rows (lemma, form, features, morpheme segmentation); derivation rows (source and target word, their parts of speech, the affix, and its type)~\cite{morphynetproject2026}.}
\resfield{What it does not record}{Senses or pronunciation: its row formats have no such fields~\cite{morphynetproject2026}.}

\resource{SPECIALIST Lexicon (National Library of Medicine)}
\resfield{What it is}{A syntactic lexicon of biomedical and general English, designed for the UMLS SPECIALIST Natural Language Processing System and released annually since 1994~\cite{nationallibrary2026}.}
\resfield{Size, as stated}{2027 release: 534,363 lexical items and 1,007,936 forms, 51.55\% of the forms multiword~\cite{nationallibrary2026a}.}
\resfield{License, as stated}{NLM's terms, not a Creative Commons license: no charge; redistributions keep the conditions; a publication must name the release number and date and describe any modifications; no proprietary rights may be claimed~\cite{nationallibrary2026b}.}
\resfield{What it records per entry}{Lexical records from which its tools retrieve syntactic category, inflectional and spelling variants, abbreviations, acronyms, and derivational variants; separate files of derivation pairs, synonyms, and antonyms~\cite{nationallibrary2026}.}
\resfield{What it does not record}{Not stated in the pages read.}

\resource{CMU Pronouncing Dictionary (CMUdict) 0.7b}
\resfield{What it is}{A machine-readable pronunciation dictionary for North American English, maintained by Carnegie Mellon University's Speech Group~\cite{carnegiemellon2026,carnegiemellon2026a}.}
\resfield{Size, as stated}{"Over 134,000 words and their pronunciations"~\cite{carnegiemellon2026}.}
\resfield{License, as stated}{A license file, Copyright 1993–2015 Carnegie Mellon University, that permits redistribution and use with or without modification if the copyright notice and disclaimer are kept~\cite{carnegiemellon2026a}; the README calls research and commercial use completely unrestricted and asks users to acknowledge its origin~\cite{carnegiemellon2026a}.}
\resfield{What it records per entry}{A pronunciation in the ARPAbet set of 39 phonemes, with a lexical stress marker on each vowel~\cite{carnegiemellon2026}.}
\resfield{What it does not record}{Anything but the word and its pronunciation: its entries map words to pronunciations~\cite{carnegiemellon2026}.}

\subsection*{Word lists}

\resource{SCOWL (SCOWLv1), release 2020.12.07. Partly game-derived at sizes 80 and above (see what it does not record).}
\resfield{What it is}{Spell Checking Oriented Word Lists: word lists split by size and spelling category, intended for use in spell checkers~\cite{atkinson2024}. As of February 17, 2026, the original format is no longer maintained~\cite{atkinson2026}.}
\resfield{Size, as stated}{The SCOWLv1 README (commit of July 22, 2024) counts 658,213 entries through size 95 in the American and English categories, names included~\cite{atkinson2024}.}
\resfield{License, as stated}{Copyright 2000–2018 Kevin Atkinson: permission to use, copy, modify, distribute, and sell, provided the notice appears in all copies; the notices of its components also apply, among them UKACD's, which requires that its notice be displayed and its text included verbatim on redistribution~\cite{atkinson2020}.}
\resfield{What it records per entry}{Membership by spelling category (American, British, Canadian, Australian, and variant levels), sub-category (abbreviations, contractions, proper names, upper case, words), and size (10 to 95)~\cite{atkinson2024}.}
\resfield{What it does not record}{Hyphenated and open (spaced) compounds, which the original version excluded~\cite{atkinson2026a}. Note: size 80 includes all the lists of the ENABLE supplement package except one~\cite{atkinson2020}, and SCOWL's successor documents that package's lists of additional and stale OSPD words~\cite{atkinson2026b}.}

\resource{English Speller Database (ESDB, formerly SCOWLv2). Partly game-derived at sizes 80 and 85 (see what it does not record).}
\resfield{What it is}{SCOWL's successor: a SQLite database and text file on how common a word is, dialect spelling differences, variants, and basic part of speech and inflection~\cite{atkinson2026a}. It is a work in progress, and the database itself will not have an official release until it stabilizes~\cite{atkinson2026c}.}
\resfield{Size, as stated}{No entry count stated; the minimum size is 35, and size 95 is not included~\cite{atkinson2026a}.}
\resfield{License, as stated}{Copyright 2000–2026 Kevin Atkinson, the same notice-retention grant, which the README calls "MIT-like"; further notices apply to Australian output, to generated lists larger than size 80 (UKACD), and possibly WordNet~\cite{atkinson2026a,atkinson2026d}.}
\resfield{What it records per entry}{Size, region, spelling and variant level, a lemma with its part of speech, inflected forms, definition notes, usage notes (among them offensive and vulgar), and sometimes a tag naming the source list~\cite{atkinson2026a}. Unlike SCOWLv1, it includes hyphenated and open compounds~\cite{atkinson2026a}.}
\resfield{What it does not record}{Note: its sources document lists additional OSPD words, added at size 80, and stale OSPD words, at size 85, both from the ENABLE2K supplement files~\cite{atkinson2026b}.}

\resource{LibreOffice hunspell dictionaries en\_US, en\_CA, en\_AU, en\_GB}
\resfield{What it is}{Spell-checker dictionaries. The en\_US, en\_CA, and en\_AU README (version 2020.12.07) says these dictionaries are made by a SCOWL script and that the normal ones correspond to SCOWL size 60; en\_GB began from a subset of Kevin Atkinson's list and has been extensively updated by others~\cite{libreofficecontributors2026}.}
\resfield{Size, as stated}{No entry count stated in the READMEs~\cite{libreofficecontributors2026}.}
\resfield{License, as stated}{en\_US, en\_CA, en\_AU: they come directly from SCOWL and are under SCOWL's copyright, with an affix file under Geoff Kuenning's BSD license; en\_GB: "covered by his original LGPL licence", with the affix file also under the LGPL~\cite{libreofficecontributors2026}.}
\resfield{What it records per entry}{Words and affix rules that generate inflected forms, 1 spelling variety per dictionary; en\_GB adds hyphenated compounds, place names, possessives, plurals, and a morphological extension based on Atkinson's part-of-speech and inflection data~\cite{libreofficecontributors2026}.}
\resfield{What it does not record}{Most alternative spellings: the size-60 dictionaries generally include 1 variant per word~\cite{libreofficecontributors2026}.}

\resource{ENABLE (ENABLE2K). A word-game list; it states it is not the tournament list.}
\resfield{What it is}{The Enhanced North American Benchmark Lexicon, a word list for word games researched and compiled by Alan Beale. Its README says the master list is not the proprietary TWL98, and that it reflects the OSPD changes in effect from January 1, 1998~\cite{cooper2026}.}
\resfield{Size, as stated}{173,528 words in the Y2K update; the first edition had 100,701~\cite{cooper2026}.}
\resfield{License, as stated}{The master list is "formally released into the Public Domain"; credit to the originators is requested~\cite{cooper2026}.}
\resfield{What it records per entry}{Membership in a plain ASCII word list~\cite{cooper2026}.}
\resfield{What it does not record}{Anything beyond the string~\cite{cooper2026}.}

\resource{12dicts 6.0.2 (June 2016). Partly game-derived (some inflections in 2of12inf).}
\resfield{What it is}{Alan Beale's word lists of common words; the classic lists record agreement among 12 source dictionaries of American English, 8 for learners of English and 4 desk dictionaries~\cite{beale2026}.}
\resfield{Size, as stated}{Classic lists from about 22,000 to about 82,000 words~\cite{beale2026}.}
\resfield{License, as stated}{Released to the public domain, with acknowledgment requested; 2of12inf and the 2+2+3 lists depend on AGID, which keeps them out of the public domain~\cite{beale2026}.}
\resfield{What it records per entry}{Membership in lists defined by how many source dictionaries hold a word; 2of12inf adds inflections, marking the plurals of uncountable nouns~\cite{beale2026}.}
\resfield{What it does not record}{Note: the README says some 2of12inf inflections came from Scrabble-related sources, and that no later list with inflections includes those plurals~\cite{beale2026}.}

\resource{Moby Words II (Grady Ward). Game-derived: the 2 crossword files.}
\resfield{What it is}{Word lists by Grady Ward, released on Project Gutenberg on May 1, 2002~\cite{ward2002}.}
\resfield{Size, as stated}{Separate lists, among them 354,984 single words, 256,772 compound words, 74,550 common dictionary words, 21,986 names, 10,196 places, and 113,809 crossword-game words~\cite{ward2002}.}
\resfield{License, as stated}{"Public Domain material by grant from the author, January, 2001"~\cite{ward2002}.}
\resfield{What it records per entry}{Membership in each list~\cite{ward2002}.}
\resfield{What it does not record}{Accents, which were stripped. The crossword file is described as compatible with the first edition of the OSPD, and a delta file makes it compatible with the second~\cite{ward2002}.}

\resource{Moby Part-of-Speech II (Grady Ward)}
\resfield{What it is}{A part-of-speech list, the second edition, with about 15,000 new entries~\cite{ward2026}.}
\resfield{Size, as stated}{No count stated in its documentation~\cite{ward2026}.}
\resfield{License, as stated}{The same 2001 public-domain grant~\cite{ward2026}.}
\resfield{What it records per entry}{A word or phrase with 1 or more of 15 part-of-speech codes, the principal use first~\cite{ward2026}.}
\resfield{What it does not record}{Many -ed, -ing, -ly, and -ic forms, which are not listed explicitly~\cite{ward2026}.}

\resource{YAWL (Yet Another Word List) 0.3.2.03. Game-derived.}
\resfield{What it is}{A word-game list by M. Leo Cooper, a superset of ENABLE, OSW, and lists by Alan Beale; rehosted on GitHub by Aaron Bull Schaefer~\cite{cooper2026a}.}
\resfield{Size, as stated}{264,097 words~\cite{cooper2026a}.}
\resfield{License, as stated}{The list "is in the Public Domain"~\cite{cooper2026a}.}
\resfield{What it records per entry}{Membership~\cite{cooper2026a}.}
\resfield{What it does not record}{Anything beyond the string~\cite{cooper2026a}.}

\resource{Letterpress word list. Game-derived.}
\resfield{What it is}{The word list behind the Letterpress game, loosely based on other word lists; its rule is that a word allowed in another popular word game such as Scrabble should be allowed, and it names YAWL, ENABLE2K, and CSW12 as useful lists~\cite{lorenbrichterwordsmaintainers2026}.}
\resfield{Size, as stated}{Not stated~\cite{lorenbrichterwordsmaintainers2026}.}
\resfield{License, as stated}{CC0 1.0 Universal~\cite{lorenbrichterwordsmaintainers2026}.}
\resfield{What it records per entry}{Membership in 1 sorted file per language; US, UK, and Australian spellings combined~\cite{lorenbrichterwordsmaintainers2026}.}
\resfield{What it does not record}{Proper nouns and acronyms, excluded by rule with a few kept for historical reasons; which spelling variety an entry belongs to~\cite{lorenbrichterwordsmaintainers2026}.}

\resource{Wordnik wordlist. A word-game list; the source of its selection is not stated.}
\resfield{What it is}{A list of English words commonly used in word games, for game developers~\cite{wordnik2026a}. Wordnik's paid game dataset, which the README offers to those who need more data, states that it takes its playable words from ENABLE and the Tournament Word List~\cite{wordnik2026}.}
\resfield{Size, as stated}{Not stated~\cite{wordnik2026a}.}
\resfield{License, as stated}{MIT License, Copyright 2020 Wordnik~\cite{wordnik2026a}.}
\resfield{What it records per entry}{Membership~\cite{wordnik2026a}.}
\resfield{What it does not record}{Definitions, parts of speech, frequency, and offensive-term labels, which the README offers in a licensed dataset~\cite{wordnik2026a}.}

\resource{dwyl/english-words}
\resfield{What it is}{A text file of English words extracted from a word list once published by Infochimps~\cite{dwyl2026}.}
\resfield{Size, as stated}{"Over 466k" words in the README; 479k in the repository's description~\cite{dwyl2026}.}
\resfield{License, as stated}{The Unlicense, beside a README that says the copyright still belongs to the original publisher~\cite{dwyl2026}.}
\resfield{What it records per entry}{1 string per line; a second file keeps only letters-only strings~\cite{dwyl2026}.}
\resfield{What it does not record}{Anything beyond the string~\cite{dwyl2026}.}

\resource{Webster's Second word lists (web2, web2a)}
\resfield{What it is}{Word lists in FreeBSD's \texttt{share/dict} directory, introduced as Webster's Second International; web2a holds hyphenated terms and noun and adverbial phrases. FreeBSD maintains web2 as a list of current American English spellings and has removed a few entries~\cite{woods2026}.}
\resfield{Size, as stated}{web2: 234,936 words; no count for web2a~\cite{woods2026}.}
\resfield{License, as stated}{No license stated; the README says the 1934 copyright "has lapsed, according to the supplier"~\cite{woods2026}. The Catalog of Copyright Entries records renewals of the dictionary's 1934 copyrights on December 18, 1961~\cite{libraryof1961}.}
\resfield{What it records per entry}{Membership~\cite{woods2026}.}
\resfield{What it does not record}{Anything beyond the string~\cite{woods2026}.}

\subsection*{Frequency resources and norms}

\resource{wordfreq 3.1.1 (Robyn Speer)}
\resfield{What it is}{A Python library of word frequencies in over 40 languages, combining several sources; its data is a snapshot of usage through about 2021 and is unlikely to be updated again~\cite{speer2026,speer2026a}. 3.1.1, uploaded November 21, 2023, is the newest release on PyPI~\cite{pythonpackage2026}.}
\resfield{Size, as stated}{No English word count stated; its table lists 7 sources for English~\cite{speer2026}.}
\resfield{License, as stated}{Code under the Apache license; data files that may be redistributed under CC BY-SA 4.0~\cite{speer2026}. The current README adds that converting the data to CSV does not follow the license, because CSV has no place for attribution~\cite{speer2026}.}
\resfield{What it records per entry}{A frequency per word, also on the Zipf scale (the base-10 logarithm of occurrences per billion words)~\cite{speer2026}.}
\resfield{What it does not record}{Words below its cutoff: the lowest Zipf value in the large lists is 1.0, and a word not in the list returns a default value of 0, not a missing value~\cite{speer2026}.}

\resource{Google Books Ngram, English, version 20200217}
\resfield{What it is}{N-gram counts from the books Google has scanned; the downloadable sets were generated in February 2020 (version 3), July 2012 (version 2), and July 2009 (version 1)~\cite{googlebooks2026}.}
\resfield{Size, as stated}{No English 1-gram type count stated~\cite{googlebooks2026}. The first corpus held 5,195,769 books, with 361 billion English words~\cite{michel2011}; the 2012 edition held 8,116,746 books, the English part close to half a trillion words~\cite{lin2012}.}
\resfield{License, as stated}{"This compilation is licensed under a Creative Commons Attribution 3.0 Unported License"~\cite{googlebooks2026}; the Viewer's page says its graphs and data may be freely used for any purpose, with acknowledgment appreciated~\cite{googlebooks2026a}.}
\resfield{What it records per entry}{Per n-gram and year, a match count and a volume count (the version 2 format)~\cite{googlebooks2026}.}
\resfield{What it does not record}{N-grams that occur 40 times or fewer across the corpus, which versions 1 and 2 omit~\cite{googlebooks2026}; how often a book was read, since the corpus holds one copy of each book~\cite{pechenick2015}.}

\resource{SUBTLEX-US}
\resfield{What it is}{Word frequencies from American film and television subtitles~\cite{brysbaert2009}.}
\resfield{Size, as stated}{74,286 words in the corpus, 60,384 of them with a frequency above 1~\cite{ghentuniversity2026}, from 51.0 million words in 8,388 films and episodes~\cite{brysbaert2009}.}
\resfield{License, as stated}{The 2009 paper: freely available for research purposes~\cite{brysbaert2009}. The author's copy on OSF carries a license file dated May 7, 2024 stating CC BY-NC-SA 4.0~\cite{brysbaert2024}. The Ghent University download page states no license~\cite{ghentuniversity2026}.}
\resfield{What it records per entry}{A frequency count, the number of films a word occurs in, a count of lowercase occurrences, and per-million values~\cite{ghentuniversity2026}.}
\resfield{What it does not record}{Not stated in the pages read.}

\resource{SUBTLEX-UK}
\resfield{What it is}{British English word frequencies from BBC subtitles broadcast from January 2010 to December 2012~\cite{vanheuven2014}.}
\resfield{Size, as stated}{332,987 word types in 201.7 million tokens; the main file holds 160,022 types (159,235 single words and 787 hyphenated)~\cite{vanheuven2014}.}
\resfield{License, as stated}{The author's OSF copy carries a license file dated May 7, 2024 stating CC BY-NC-SA 4.0~\cite{brysbaert2024a}.}
\resfield{What it records per entry}{Frequency, Zipf value, contextual diversity, the dominant part of speech and its lemma, and all parts of speech observed~\cite{vanheuven2014}.}
\resfield{What it does not record}{Types seen in a single broadcast, and types with digits or other non-letter symbols, which the main file leaves out~\cite{vanheuven2014}.}

\resource{English Lexicon Project}
\resfield{What it is}{Behavioral and descriptive data for lexical decision and naming studies~\cite{balota2007}.}
\resfield{Size, as stated}{40,481 words and 40,481 nonwords~\cite{balota2007}.}
\resfield{License, as stated}{No terms on the home page (archived October 14, 2026); an About page archived in 2019 says word lists from the site are for non-commercial research only and may not be used to develop speech technology~\cite{englishlexicon2026}. The live site's certificate had expired when read on October 26, 2026.}
\resfield{What it records per entry}{Lexical decision and naming responses (816 and 444 participants) and descriptive lexical statistics~\cite{balota2007}.}
\resfield{What it does not record}{Not stated in the pages read.}

\resource{Web 1T 5-gram Version 1 (LDC2006T13). Not open.}
\resfield{What it is}{English word n-grams and their counts from about 1 trillion tokens of text from publicly accessible web pages, contributed by Google~\cite{brants2006}.}
\resfield{Size, as stated}{13,588,391 unigrams; 1,024,908,267,229 tokens~\cite{brants2006}.}
\resfield{License, as stated}{The Web 1T 5-gram Version 1 Agreement of the Linguistic Data Consortium; fees apply~\cite{brants2006}.}
\resfield{What it records per entry}{N-grams of 1 to 5 words with observed frequency counts~\cite{brants2006}.}
\resfield{What it does not record}{Not stated in the pages read.}


\end{document}
